%% kapa: exam=1 round=33 year=2022 period=1 form=A subject=civil_law subject_ko=민법 seq=36 number=36
%% source: Q-Net 감정평가사 기출문제 (cst00309) artlSeq=5209997 file=raw/1008/5209997/2022년도 제33회 감정평가사 1차시험 1교시(민법,경제학,부동산학).pdf
%% conversion: kordoc 4.24.51 -> tools/gen1.py (auto); review=proofread-sonnet; flags=none
\begin{question}{36}
민사동산질권에 관한 설명으로 옳지 않은 것은?
\begin{choices}
  \item 질권자는 피담보채권의 변제를 받기 위하여 질물을 경매할 수 있고, 그 매각대금으로부터 일반채권자와 동일한 순위로 변제받는다.
  \item 질권은 양도할 수 없는 물건을 목적으로 하지 못한다.
  \item 질권은 다른 약정이 없는 한 원본, 이자, 위약금, 질권실행의 비용, 질물보존의 비용 및 채무불이행 또는 질물의 하자로 인한 손해배상의 채권을 담보한다.
  \item 질권자는 피담보채권의 변제를 받을 때까지 질물을 유치할 수 있으나 자기보다 우선권이 있는 채권자에게 대항하지 못한다.
  \item 수개의 채권을 담보하기 위하여 동일한 동산에 수개의 질권을 설정한 때에는 그 순위는 설정의 선후에 의한다.
\end{choices}
\end{question}
\answer{1}
